\documentclass[a4paper,11pt]{article}
\usepackage[utf8]{inputenc}
\usepackage[T1]{fontenc}
\usepackage[es-ES]{babel}
\usepackage[margin=2.5cm]{geometry}
\usepackage{amsmath}
\usepackage{enumitem}
\usepackage{xcolor}
\usepackage{tcolorbox}
\usepackage{graphicx}
\usepackage{hyperref}
\usepackage{listings}
\usepackage{booktabs}

% Colores
\definecolor{unidadcolor}{RGB}{30, 60, 120}
\definecolor{objetivocolor}{RGB}{0, 100, 0}
\definecolor{ejemplocolor}{RGB}{200, 200, 200}
\definecolor{advertenciacolor}{RGB}{200, 100, 100}

% Configuración de listings
\lstset{
    basicstyle=\ttfamily\small,
    keywordstyle=\color{unidadcolor}\bfseries,
    commentstyle=\color{gray}\itshape,
    stringstyle=\color{objetivocolor},
    showstringspaces=false,
    breaklines=true,
    frame=single,
    backgroundcolor=\color{gray!10}
}

\begin{document}

\begin{center}
    {\Large \textbf{Métodos de Análisis de Datos 1}}\\[0.3em]
    {\large Unidad 2: Exploración y Visualización de Datos}\\[0.2em]
    {\normalsize Clases 8 a 13}
\end{center}

\vspace{1em}

\section*{Objetivos de la unidad}

Al finalizar esta unidad, los estudiantes serán capaces de:
\begin{itemize}[leftmargin=*]
    \item Calcular e interpretar medidas de tendencia central y dispersión
    \item Analizar la distribución de variables
    \item Detectar e interpretar outliers
    \item Realizar análisis de relaciones entre variables
    \item Crear visualizaciones efectivas con matplotlib, seaborn y ggplot2
    \item Comunicar hallazgos de forma clara y precisa
\end{itemize}

\section{Clase 8: Estadística Descriptiva Univariada}

\subsection{Introducción}

El análisis exploratorio de datos (EDA - Exploratory Data Analysis) es el proceso de examinar datos para descubrir patrones, detectar anomalías, probar hipótesis y verificar supuestos mediante resúmenes estadísticos y representaciones gráficas.

\begin{tcolorbox}[colback=ejemplocolor!20, colframe=unidadcolor, title=¿Por qué es importante el EDA?]
\begin{itemize}
    \item Revela la estructura subyacente de los datos
    \item Identifica variables importantes
    \item Detecta outliers y anomalías
    \item Valida supuestos para modelos
    \item Sugiere transformaciones apropiadas
    \item Guía la selección de modelos
\end{itemize}
\end{tcolorbox}

\subsection{Tipos de variables}

\subsubsection{Variables cuantitativas}

\begin{itemize}
    \item \textbf{Discretas:} Valores contables (número de hijos, cantidad de ventas)
    \item \textbf{Continuas:} Cualquier valor en un intervalo (altura, peso, temperatura)
\end{itemize}

\subsubsection{Variables cualitativas (categóricas)}

\begin{itemize}
    \item \textbf{Nominales:} Sin orden natural (color, ciudad, género)
    \item \textbf{Ordinales:} Con orden natural (nivel educativo, satisfacción)
\end{itemize}

\subsection{Medidas de tendencia central}

Las medidas de tendencia central indican el valor "típico" o "central" de una distribución.

\subsubsection{Media aritmética}

$$\bar{x} = \frac{1}{n}\sum_{i=1}^{n} x_i = \frac{x_1 + x_2 + \cdots + x_n}{n}$$

\begin{itemize}
    \item \textbf{Ventajas:} Usa toda la información, única
    \item \textbf{Desventajas:} Sensible a outliers
    \item \textbf{Usar cuando:} Distribución simétrica, sin outliers extremos
\end{itemize}

\textbf{Python:}
\begin{lstlisting}[language=Python]
media = df['columna'].mean()
\end{lstlisting}

\textbf{R:}
\begin{lstlisting}[language=R]
media <- mean(df$columna)
\end{lstlisting}

\subsubsection{Mediana}

El valor que divide la distribución en dos mitades iguales. Para datos ordenados $x_1, x_2, \ldots, x_n$:

$$\text{Mediana} = \begin{cases}
x_{(n+1)/2} & \text{si } n \text{ es impar}\\
\frac{x_{n/2} + x_{n/2+1}}{2} & \text{si } n \text{ es par}
\end{cases}$$

\begin{itemize}
    \item \textbf{Ventajas:} Robusta a outliers
    \item \textbf{Desventajas:} No usa toda la información
    \item \textbf{Usar cuando:} Distribución asimétrica o presencia de outliers
\end{itemize}

\textbf{Python:}
\begin{lstlisting}[language=Python]
mediana = df['columna'].median()
\end{lstlisting}

\textbf{R:}
\begin{lstlisting}[language=R]
mediana <- median(df$columna)
\end{lstlisting}

\subsubsection{Moda}

El valor más frecuente en la distribución.

\begin{itemize}
    \item \textbf{Ventajas:} Aplicable a datos cualitativos
    \item \textbf{Desventajas:} Puede no ser única
    \item \textbf{Usar cuando:} Variables categóricas
\end{itemize}

\textbf{Python:}
\begin{lstlisting}[language=Python]
moda = df['columna'].mode()[0]  # Primera moda si hay múltiples
\end{lstlisting}

\textbf{R:}
\begin{lstlisting}[language=R]
# Función personalizada
moda <- function(x) {
  ux <- unique(x)
  ux[which.max(tabulate(match(x, ux)))]
}
\end{lstlisting}

\begin{tcolorbox}[colback=ejemplocolor!20, colframe=unidadcolor, title=Ejemplo: Media vs Mediana]
Salarios: \$30k, \$32k, \$35k, \$38k, \$40k, \$200k

\textbf{Media:} \$62.5k (distorsionada por el outlier)

\textbf{Mediana:} \$36.5k (valor más representativo)

En este caso, la mediana es más apropiada.
\end{tcolorbox}

\subsection{Medidas de dispersión}

Las medidas de dispersión cuantifican la variabilidad de los datos.

\subsubsection{Rango}

$$\text{Rango} = \text{Máximo} - \text{Mínimo}$$

\begin{itemize}
    \item Simple de calcular
    \item Muy sensible a outliers
    \item No considera la distribución completa
\end{itemize}

\subsubsection{Rango intercuartílico (IQR)}

$$\text{IQR} = Q_3 - Q_1$$

Donde $Q_1$ es el primer cuartil (percentil 25) y $Q_3$ es el tercer cuartil (percentil 75).

\begin{itemize}
    \item Robusto a outliers
    \item Captura la dispersión del 50\% central de los datos
\end{itemize}

\textbf{Python:}
\begin{lstlisting}[language=Python]
Q1 = df['columna'].quantile(0.25)
Q3 = df['columna'].quantile(0.75)
IQR = Q3 - Q1
\end{lstlisting}

\textbf{R:}
\begin{lstlisting}[language=R]
IQR_val <- IQR(df$columna)
\end{lstlisting}

\subsubsection{Varianza}

$$s^2 = \frac{1}{n-1}\sum_{i=1}^{n}(x_i - \bar{x})^2$$

La varianza mide el promedio de las desviaciones al cuadrado respecto a la media.

\begin{itemize}
    \item Usa toda la información
    \item Unidades al cuadrado (difícil interpretación)
    \item Sensible a outliers
\end{itemize}

\subsubsection{Desviación estándar}

$$s = \sqrt{s^2} = \sqrt{\frac{1}{n-1}\sum_{i=1}^{n}(x_i - \bar{x})^2}$$

La raíz cuadrada de la varianza, en las mismas unidades que los datos originales.

\textbf{Python:}
\begin{lstlisting}[language=Python]
varianza = df['columna'].var()
desv_std = df['columna'].std()
\end{lstlisting}

\textbf{R:}
\begin{lstlisting}[language=R]
varianza <- var(df$columna)
desv_std <- sd(df$columna)
\end{lstlisting}

\begin{tcolorbox}[colback=ejemplocolor!20, colframe=unidadcolor, title=Interpretación de la desviación estándar]
Si los datos siguen aproximadamente una distribución normal:
\begin{itemize}
    \item 68\% de los datos están dentro de $\pm 1$ desviación estándar de la media
    \item 95\% están dentro de $\pm 2$ desviaciones estándar
    \item 99.7\% están dentro de $\pm 3$ desviaciones estándar
\end{itemize}

Esta es la \textbf{regla empírica} o \textbf{regla 68-95-99.7}.
\end{tcolorbox}

\subsubsection{Coeficiente de variación}

$$CV = \frac{s}{\bar{x}} \times 100\%$$

Medida de dispersión relativa, útil para comparar la variabilidad de variables con diferentes unidades o escalas.

\begin{itemize}
    \item Adimensional (porcentaje)
    \item Permite comparar dispersión entre diferentes variables
    \item Solo para variables con razón (no para temperaturas en Celsius, por ejemplo)
\end{itemize}

\textbf{Python:}
\begin{lstlisting}[language=Python]
cv = (df['columna'].std() / df['columna'].mean()) * 100
\end{lstlisting}

\textbf{R:}
\begin{lstlisting}[language=R]
cv <- (sd(df$columna) / mean(df$columna)) * 100
\end{lstlisting}

\subsection{Medidas de forma}

\subsubsection{Asimetría (Skewness)}

Mide la falta de simetría de la distribución.

$$\text{Skewness} = \frac{1}{n}\sum_{i=1}^{n}\left(\frac{x_i - \bar{x}}{s}\right)^3$$

\begin{itemize}
    \item \textbf{Skewness = 0:} Distribución simétrica
    \item \textbf{Skewness > 0:} Asimetría positiva (cola derecha más larga)
    \item \textbf{Skewness < 0:} Asimetría negativa (cola izquierda más larga)
\end{itemize}

\textbf{Python:}
\begin{lstlisting}[language=Python]
from scipy import stats
asimetria = stats.skew(df['columna'])
# O con pandas
asimetria = df['columna'].skew()
\end{lstlisting}

\textbf{R:}
\begin{lstlisting}[language=R]
library(moments)
asimetria <- skewness(df$columna)
\end{lstlisting}

\subsubsection{Curtosis (Kurtosis)}

Mide el "apuntamiento" o concentración de valores alrededor de la media.

$$\text{Kurtosis} = \frac{1}{n}\sum_{i=1}^{n}\left(\frac{x_i - \bar{x}}{s}\right)^4 - 3$$

\begin{itemize}
    \item \textbf{Kurtosis = 0:} Mesocúrtica (normal)
    \item \textbf{Kurtosis > 0:} Leptocúrtica (colas pesadas, pico agudo)
    \item \textbf{Kurtosis < 0:} Platicúrtica (colas ligeras, pico plano)
\end{itemize}

\textbf{Python:}
\begin{lstlisting}[language=Python]
from scipy import stats
curtosis = stats.kurtosis(df['columna'])
# O con pandas
curtosis = df['columna'].kurtosis()
\end{lstlisting}

\textbf{R:}
\begin{lstlisting}[language=R]
library(moments)
curtosis <- kurtosis(df$columna) - 3  # Exceso de curtosis
\end{lstlisting}

\section{Clase 9: Análisis de Distribuciones}

\subsection{Visualización de distribuciones univariadas}

\subsubsection{Histogramas}

Los histogramas muestran la distribución de una variable continua dividiendo el rango en intervalos (bins) y contando cuántas observaciones caen en cada intervalo.

\textbf{Python:}
\begin{lstlisting}[language=Python]
import matplotlib.pyplot as plt

plt.hist(df['columna'], bins=30, edgecolor='black')
plt.xlabel('Variable')
plt.ylabel('Frecuencia')
plt.title('Histograma')
plt.show()
\end{lstlisting}

\textbf{R:}
\begin{lstlisting}[language=R]
library(ggplot2)

ggplot(df, aes(x = columna)) +
  geom_histogram(bins = 30, color = "black", fill = "steelblue") +
  labs(x = "Variable", y = "Frecuencia", title = "Histograma")
\end{lstlisting}

\begin{tcolorbox}[colback=advertenciacolor!20, colframe=red, title=Cuidado con el número de bins]
\begin{itemize}
    \item Muy pocos bins: Se pierde información de la forma
    \item Demasiados bins: Ruido, difícil interpretación
    \item \textbf{Regla de Sturges:} $k = 1 + \log_2(n)$
    \item \textbf{Regla de Rice:} $k = 2n^{1/3}$
    \item Lo mejor: Experimentar con diferentes valores
\end{itemize}
\end{tcolorbox}

\subsubsection{Gráficos de densidad}

Estimación suave de la función de densidad de probabilidad.

\textbf{Python:}
\begin{lstlisting}[language=Python]
import seaborn as sns

sns.kdeplot(data=df, x='columna', fill=True)
plt.xlabel('Variable')
plt.ylabel('Densidad')
plt.title('Gráfico de Densidad')
plt.show()
\end{lstlisting}

\textbf{R:}
\begin{lstlisting}[language=R]
ggplot(df, aes(x = columna)) +
  geom_density(fill = "steelblue", alpha = 0.5) +
  labs(x = "Variable", y = "Densidad", title = "Gráfico de Densidad")
\end{lstlisting}

\subsubsection{Boxplots (Diagramas de caja)}

Resumen visual de los cinco números: mínimo, Q1, mediana, Q3, máximo.

\textbf{Componentes:}
\begin{itemize}
    \item \textbf{Caja:} Rango intercuartílico (IQR)
    \item \textbf{Línea central:} Mediana
    \item \textbf{Bigotes:} Se extienden hasta 1.5 × IQR desde los cuartiles
    \item \textbf{Puntos:} Outliers (fuera de los bigotes)
\end{itemize}

\textbf{Python:}
\begin{lstlisting}[language=Python]
plt.boxplot(df['columna'])
plt.ylabel('Valor')
plt.title('Boxplot')
plt.show()

# O con seaborn
sns.boxplot(data=df, y='columna')
plt.show()
\end{lstlisting}

\textbf{R:}
\begin{lstlisting}[language=R]
ggplot(df, aes(y = columna)) +
  geom_boxplot(fill = "steelblue") +
  labs(y = "Valor", title = "Boxplot")
\end{lstlisting}

\subsubsection{Violin plots}

Combinación de boxplot y gráfico de densidad.

\textbf{Python:}
\begin{lstlisting}[language=Python]
sns.violinplot(data=df, y='columna')
plt.show()
\end{lstlisting}

\textbf{R:}
\begin{lstlisting}[language=R]
ggplot(df, aes(y = columna)) +
  geom_violin(fill = "steelblue") +
  labs(y = "Valor", title = "Violin Plot")
\end{lstlisting}

\subsection{Pruebas de normalidad}

Muchos métodos estadísticos asumen que los datos siguen una distribución normal.

\subsubsection{Q-Q Plot (Quantile-Quantile Plot)}

Compara los cuantiles de los datos con los cuantiles teóricos de una distribución normal.

\begin{itemize}
    \item Si los puntos siguen la línea diagonal, los datos son aproximadamente normales
    \item Desviaciones en los extremos indican colas pesadas o ligeras
\end{itemize}

\textbf{Python:}
\begin{lstlisting}[language=Python]
from scipy import stats
import matplotlib.pyplot as plt

stats.probplot(df['columna'], dist="norm", plot=plt)
plt.title('Q-Q Plot')
plt.show()
\end{lstlisting}

\textbf{R:}
\begin{lstlisting}[language=R]
qqnorm(df$columna)
qqline(df$columna)
\end{lstlisting}

\subsubsection{Test de Shapiro-Wilk}

Prueba estadística de normalidad.

\begin{itemize}
    \item $H_0$: Los datos provienen de una distribución normal
    \item $H_1$: Los datos no provienen de una distribución normal
    \item Si $p < 0.05$, rechazar $H_0$ (datos no normales)
\end{itemize}

\textbf{Python:}
\begin{lstlisting}[language=Python]
from scipy import stats

statistic, p_value = stats.shapiro(df['columna'])
print(f"Estadístico: {statistic}, p-value: {p_value}")

if p_value > 0.05:
    print("Los datos parecen normales (no rechazamos H0)")
else:
    print("Los datos NO parecen normales (rechazamos H0)")
\end{lstlisting}

\textbf{R:}
\begin{lstlisting}[language=R]
resultado <- shapiro.test(df$columna)
print(resultado)
\end{lstlisting}

\begin{tcolorbox}[colback=advertenciacolor!20, colframe=red, title=Limitaciones del test de Shapiro-Wilk]
\begin{itemize}
    \item Muy sensible con muestras grandes (detecta desviaciones triviales)
    \item Funciona mejor con muestras de tamaño moderado (50-2000)
    \item Combinar con inspección visual (Q-Q plot, histograma)
\end{itemize}
\end{tcolorbox}

\subsection{Transformaciones para normalizar}

Cuando los datos no son normales, se pueden aplicar transformaciones:

\subsubsection{Transformación logarítmica}

$$y = \log(x)$$

\begin{itemize}
    \item Útil para datos con asimetría positiva
    \item Requiere $x > 0$
    \item Si hay ceros: $y = \log(x + 1)$
\end{itemize}

\subsubsection{Transformación de raíz cuadrada}

$$y = \sqrt{x}$$

\begin{itemize}
    \item Menos drástica que log
    \item Útil para datos de conteo
\end{itemize}

\subsubsection{Transformación Box-Cox}

$$y = \begin{cases}
\frac{x^\lambda - 1}{\lambda} & \text{si } \lambda \neq 0\\
\log(x) & \text{si } \lambda = 0
\end{cases}$$

Encuentra automáticamente la mejor transformación potencia.

\textbf{Python:}
\begin{lstlisting}[language=Python]
from scipy import stats

# Box-Cox (requiere datos positivos)
transformed, lambda_param = stats.boxcox(df['columna'])
print(f"Lambda óptimo: {lambda_param}")
\end{lstlisting}

\textbf{R:}
\begin{lstlisting}[language=R]
library(MASS)

# Box-Cox
bc <- boxcox(columna ~ 1, data = df)
lambda <- bc$x[which.max(bc$y)]
\end{lstlisting}

\section{Clase 10: Detección de Outliers}

\subsection{¿Qué son los outliers?}

Los outliers (valores atípicos) son observaciones que se desvían significativamente del resto de los datos.

\begin{tcolorbox}[colback=ejemplocolor!20, colframe=unidadcolor, title=Tipos de outliers]
\begin{itemize}
    \item \textbf{Errores de medición:} Deben corregirse o eliminarse
    \item \textbf{Errores de entrada:} Deben corregirse
    \item \textbf{Valores genuinos pero raros:} Pueden ser los más interesantes (fraudes, casos extremos)
    \item \textbf{Muestreo de otra población:} Considerar análisis separado
\end{itemize}
\end{tcolorbox}

\subsection{Métodos univariados}

\subsubsection{Método IQR (Rango Intercuartílico)}

Define outliers como valores fuera del rango:

$$[Q_1 - 1.5 \times \text{IQR}, \quad Q_3 + 1.5 \times \text{IQR}]$$

\textbf{Python:}
\begin{lstlisting}[language=Python]
Q1 = df['columna'].quantile(0.25)
Q3 = df['columna'].quantile(0.75)
IQR = Q3 - Q1

lower_bound = Q1 - 1.5 * IQR
upper_bound = Q3 + 1.5 * IQR

outliers = df[(df['columna'] < lower_bound) | (df['columna'] > upper_bound)]
print(f"Outliers detectados: {len(outliers)}")
\end{lstlisting}

\textbf{R:}
\begin{lstlisting}[language=R]
Q1 <- quantile(df$columna, 0.25)
Q3 <- quantile(df$columna, 0.75)
IQR_val <- Q3 - Q1

lower_bound <- Q1 - 1.5 * IQR_val
upper_bound <- Q3 + 1.5 * IQR_val

outliers <- df[df$columna < lower_bound | df$columna > upper_bound, ]
cat("Outliers detectados:", nrow(outliers), "\n")
\end{lstlisting}

\subsubsection{Método Z-score}

Define outliers como valores con $|z| > 3$ (o a veces $|z| > 2.5$).

$$z = \frac{x - \bar{x}}{s}$$

\textbf{Python:}
\begin{lstlisting}[language=Python]
from scipy import stats

z_scores = np.abs(stats.zscore(df['columna']))
outliers = df[z_scores > 3]
print(f"Outliers detectados: {len(outliers)}")
\end{lstlisting}

\textbf{R:}
\begin{lstlisting}[language=R]
z_scores <- abs((df$columna - mean(df$columna)) / sd(df$columna))
outliers <- df[z_scores > 3, ]
cat("Outliers detectados:", nrow(outliers), "\n")
\end{lstlisting}

\begin{tcolorbox}[colback=advertenciacolor!20, colframe=red, title=Z-score vs IQR]
\textbf{Z-score:}
\begin{itemize}
    \item Asume distribución normal
    \item Sensible a outliers extremos (la media y std se ven afectadas)
\end{itemize}

\textbf{IQR:}
\begin{itemize}
    \item No asume distribución específica
    \item Robusto a outliers
    \item Más conservador (detecta menos outliers)
\end{itemize}
\end{tcolorbox}

\subsubsection{Método de desviación absoluta mediana (MAD)}

Más robusto que Z-score:

$$\text{MAD} = \text{mediana}(|x_i - \text{mediana}(x)|)$$

$$\text{MAD-score} = \frac{0.6745(x_i - \text{mediana}(x))}{\text{MAD}}$$

Outliers: $|\text{MAD-score}| > 3.5$

\textbf{Python:}
\begin{lstlisting}[language=Python]
from scipy import stats

mad = stats.median_abs_deviation(df['columna'])
median = df['columna'].median()
mad_scores = 0.6745 * (df['columna'] - median) / mad

outliers = df[np.abs(mad_scores) > 3.5]
\end{lstlisting}

\textbf{R:}
\begin{lstlisting}[language=R]
mad_val <- mad(df$columna)
median_val <- median(df$columna)
mad_scores <- 0.6745 * (df$columna - median_val) / mad_val

outliers <- df[abs(mad_scores) > 3.5, ]
\end{lstlisting}

\subsection{Métodos multivariados}

\subsubsection{Distancia de Mahalanobis}

Mide la distancia de un punto al centro de la distribución multivariada, considerando la covarianza.

$$D^2 = (x - \mu)^T \Sigma^{-1} (x - \mu)$$

Donde $\mu$ es el vector de medias y $\Sigma$ es la matriz de covarianza.

\textbf{Python:}
\begin{lstlisting}[language=Python]
from scipy.spatial.distance import mahalanobis

# Seleccionar variables numéricas
X = df[['var1', 'var2', 'var3']].values

# Calcular medias y covarianza
mean = X.mean(axis=0)
cov = np.cov(X.T)
cov_inv = np.linalg.inv(cov)

# Distancias de Mahalanobis
distances = [mahalanobis(x, mean, cov_inv) for x in X]

# Outliers: distancia > umbral (ej: chi-cuadrado con p=0.001)
from scipy.stats import chi2
threshold = chi2.ppf(0.999, df=X.shape[1])
outliers = df[distances > threshold]
\end{lstlisting}

\textbf{R:}
\begin{lstlisting}[language=R]
library(mvoutlier)

# Seleccionar variables
X <- df[, c('var1', 'var2', 'var3')]

# Distancia de Mahalanobis
distances <- mahalanobis(X, colMeans(X), cov(X))

# Outliers
threshold <- qchisq(0.999, df = ncol(X))
outliers <- df[distances > threshold, ]
\end{lstlisting}

\subsubsection{Isolation Forest}

Algoritmo de machine learning para detección de anomalías.

\textbf{Idea:} Los outliers son más fáciles de aislar en particiones aleatorias del espacio.

\textbf{Python:}
\begin{lstlisting}[language=Python]
from sklearn.ensemble import IsolationForest

# Entrenar modelo
iso_forest = IsolationForest(contamination=0.1, random_state=42)
predictions = iso_forest.fit_predict(df[['var1', 'var2', 'var3']])

# -1 indica outlier, 1 indica normal
outliers = df[predictions == -1]
print(f"Outliers detectados: {len(outliers)}")
\end{lstlisting}

\textbf{R:}
\begin{lstlisting}[language=R]
library(isotree)

# Entrenar modelo
modelo <- isolation.forest(df[, c('var1', 'var2', 'var3')])
scores <- predict(modelo, df[, c('var1', 'var2', 'var3')])

# Outliers: scores altos
threshold <- quantile(scores, 0.9)
outliers <- df[scores > threshold, ]
\end{lstlisting}

\subsection{Tratamiento de outliers}

\begin{enumerate}
    \item \textbf{Investigar:} ¿Son errores o valores legítimos?
    \item \textbf{Documentar:} Registrar qué se encontró y qué se hizo
    \item \textbf{Opciones:}
    \begin{itemize}
        \item \textbf{Conservar:} Si son valores genuinos de interés
        \item \textbf{Eliminar:} Si son errores claros (con justificación)
        \item \textbf{Transformar:} Aplicar transformaciones (log, sqrt)
        \item \textbf{Winsorizar:} Reemplazar con percentiles (ej: percentil 99)
        \item \textbf{Análisis separado:} Modelo con y sin outliers
        \item \textbf{Usar métodos robustos:} Que no se vean afectados por outliers
    \end{itemize}
\end{enumerate}

\section{Clase 11: Análisis Bivariado}

\subsection{Relación entre variables cuantitativas}

\subsubsection{Scatterplots (gráficos de dispersión)}

Visualizan la relación entre dos variables continuas.

\textbf{Python:}
\begin{lstlisting}[language=Python]
plt.scatter(df['var1'], df['var2'], alpha=0.5)
plt.xlabel('Variable 1')
plt.ylabel('Variable 2')
plt.title('Scatterplot')
plt.show()

# Con seaborn (más opciones)
sns.scatterplot(data=df, x='var1', y='var2', hue='categoria')
plt.show()
\end{lstlisting}

\textbf{R:}
\begin{lstlisting}[language=R]
ggplot(df, aes(x = var1, y = var2)) +
  geom_point(alpha = 0.5) +
  labs(x = "Variable 1", y = "Variable 2", title = "Scatterplot")

# Con color por categoría
ggplot(df, aes(x = var1, y = var2, color = categoria)) +
  geom_point(alpha = 0.5)
\end{lstlisting}

\subsubsection{Correlación lineal (Pearson)}

Mide la fuerza y dirección de la relación lineal entre dos variables:

$$r = \frac{\sum_{i=1}^{n}(x_i - \bar{x})(y_i - \bar{y})}{\sqrt{\sum_{i=1}^{n}(x_i - \bar{x})^2}\sqrt{\sum_{i=1}^{n}(y_i - \bar{y})^2}}$$

\begin{itemize}
    \item $r \in [-1, 1]$
    \item $r = 1$: Correlación positiva perfecta
    \item $r = -1$: Correlación negativa perfecta
    \item $r = 0$: Sin correlación lineal
    \item $|r| < 0.3$: Correlación débil
    \item $0.3 \leq |r| < 0.7$: Correlación moderada
    \item $|r| \geq 0.7$: Correlación fuerte
\end{itemize}

\textbf{Python:}
\begin{lstlisting}[language=Python]
correlacion = df['var1'].corr(df['var2'])
print(f"Correlación: {correlacion:.3f}")

# Matriz de correlación
correlation_matrix = df[['var1', 'var2', 'var3']].corr()
print(correlation_matrix)
\end{lstlisting}

\textbf{R:}
\begin{lstlisting}[language=R]
correlacion <- cor(df$var1, df$var2)
cat("Correlación:", correlacion, "\n")

# Matriz de correlación
correlation_matrix <- cor(df[, c('var1', 'var2', 'var3')])
print(correlation_matrix)
\end{lstlisting}

\begin{tcolorbox}[colback=advertenciacolor!20, colframe=red, title=Limitaciones de la correlación de Pearson]
\begin{itemize}
    \item Solo mide relaciones \textbf{lineales}
    \item Sensible a outliers
    \item Asume variables continuas y distribución bivariada normal
    \item \textbf{Correlación no implica causalidad}
\end{itemize}
\end{tcolorbox}

\subsubsection{Correlación de Spearman}

Correlación basada en rangos, más robusta a outliers y no asume linealidad.

\textbf{Python:}
\begin{lstlisting}[language=Python]
from scipy.stats import spearmanr

corr, p_value = spearmanr(df['var1'], df['var2'])
print(f"Correlación de Spearman: {corr:.3f}, p-value: {p_value:.4f}")
\end{lstlisting}

\textbf{R:}
\begin{lstlisting}[language=R]
resultado <- cor.test(df$var1, df$var2, method = "spearman")
print(resultado)
\end{lstlisting}

\subsubsection{Heatmap de correlación}

Visualización de matriz de correlación.

\textbf{Python:}
\begin{lstlisting}[language=Python]
import seaborn as sns

correlation_matrix = df.corr()

plt.figure(figsize=(10, 8))
sns.heatmap(correlation_matrix, annot=True, cmap='coolwarm', 
            center=0, fmt='.2f', square=True)
plt.title('Matriz de Correlación')
plt.show()
\end{lstlisting}

\textbf{R:}
\begin{lstlisting}[language=R]
library(corrplot)

correlation_matrix <- cor(df[, sapply(df, is.numeric)])

corrplot(correlation_matrix, method = "color", type = "upper",
         addCoef.col = "black", tl.col = "black", tl.srt = 45)
\end{lstlisting}

\subsection{Relación entre variable categórica y cuantitativa}

\subsubsection{Boxplots por grupo}

\textbf{Python:}
\begin{lstlisting}[language=Python]
sns.boxplot(data=df, x='categoria', y='valor')
plt.title('Distribución por Categoría')
plt.show()
\end{lstlisting}

\textbf{R:}
\begin{lstlisting}[language=R]
ggplot(df, aes(x = categoria, y = valor, fill = categoria)) +
  geom_boxplot() +
  labs(title = "Distribución por Categoría")
\end{lstlisting}

\subsubsection{Violin plots por grupo}

\textbf{Python:}
\begin{lstlisting}[language=Python]
sns.violinplot(data=df, x='categoria', y='valor')
plt.show()
\end{lstlisting}

\textbf{R:}
\begin{lstlisting}[language=R]
ggplot(df, aes(x = categoria, y = valor, fill = categoria)) +
  geom_violin() +
  labs(title = "Violin Plot por Categoría")
\end{lstlisting}

\subsection{Relación entre variables categóricas}

\subsubsection{Tablas de contingencia}

\textbf{Python:}
\begin{lstlisting}[language=Python]
# Tabla de frecuencias
tabla = pd.crosstab(df['categoria1'], df['categoria2'])
print(tabla)

# Tabla de proporciones
tabla_prop = pd.crosstab(df['categoria1'], df['categoria2'], normalize=True)
print(tabla_prop)
\end{lstlisting}

\textbf{R:}
\begin{lstlisting}[language=R]
# Tabla de frecuencias
tabla <- table(df$categoria1, df$categoria2)
print(tabla)

# Tabla de proporciones
tabla_prop <- prop.table(tabla)
print(tabla_prop)
\end{lstlisting}

\subsubsection{Test de Chi-cuadrado}

Prueba de independencia entre variables categóricas.

\begin{itemize}
    \item $H_0$: Las variables son independientes
    \item $H_1$: Las variables están asociadas
\end{itemize}

\textbf{Python:}
\begin{lstlisting}[language=Python]
from scipy.stats import chi2_contingency

tabla = pd.crosstab(df['categoria1'], df['categoria2'])
chi2, p_value, dof, expected = chi2_contingency(tabla)

print(f"Chi-cuadrado: {chi2:.3f}")
print(f"p-value: {p_value:.4f}")

if p_value < 0.05:
    print("Rechazamos H0: las variables están asociadas")
else:
    print("No rechazamos H0: las variables parecen independientes")
\end{lstlisting}

\textbf{R:}
\begin{lstlisting}[language=R]
resultado <- chisq.test(tabla)
print(resultado)
\end{lstlisting}

\section{Clase 12: Visualización con Python}

\subsection{Matplotlib: Fundamentos}

Matplotlib es la librería base para visualización en Python.

\subsubsection{Anatomía de una figura}

\begin{itemize}
    \item \textbf{Figure:} El contenedor principal
    \item \textbf{Axes:} El área donde se dibuja (puede haber múltiples)
    \item \textbf{Axis:} Los ejes x e y
    \item \textbf{Artist:} Elementos visuales (líneas, texto, etc.)
\end{itemize}

\textbf{Estructura básica:}
\begin{lstlisting}[language=Python]
import matplotlib.pyplot as plt

# Método pyplot (simple)
plt.plot([1, 2, 3], [1, 4, 9])
plt.xlabel('X')
plt.ylabel('Y')
plt.title('Título')
plt.show()

# Método orientado a objetos (más control)
fig, ax = plt.subplots(figsize=(10, 6))
ax.plot([1, 2, 3], [1, 4, 9])
ax.set_xlabel('X')
ax.set_ylabel('Y')
ax.set_title('Título')
plt.show()
\end{lstlisting}

\subsubsection{Múltiples subplots}

\begin{lstlisting}[language=Python]
fig, axes = plt.subplots(2, 2, figsize=(12, 10))

# axes es un array 2D
axes[0, 0].plot(x1, y1)
axes[0, 0].set_title('Gráfico 1')

axes[0, 1].scatter(x2, y2)
axes[0, 1].set_title('Gráfico 2')

axes[1, 0].hist(data, bins=30)
axes[1, 0].set_title('Gráfico 3')

axes[1, 1].boxplot(data)
axes[1, 1].set_title('Gráfico 4')

plt.tight_layout()
plt.show()
\end{lstlisting}

\subsubsection{Personalización}

\begin{lstlisting}[language=Python]
fig, ax = plt.subplots(figsize=(10, 6))

# Línea
ax.plot(x, y, 
        color='steelblue',    # Color
        linewidth=2,           # Grosor
        linestyle='--',        # Estilo: '-', '--', '-.', ':'
        marker='o',            # Marcador
        markersize=8,          # Tamaño marcador
        alpha=0.7,             # Transparencia
        label='Datos')         # Etiqueta para leyenda

# Configuración de ejes
ax.set_xlabel('X', fontsize=12)
ax.set_ylabel('Y', fontsize=12)
ax.set_title('Título', fontsize=14, fontweight='bold')

# Límites
ax.set_xlim([0, 10])
ax.set_ylim([0, 100])

# Grilla
ax.grid(True, alpha=0.3)

# Leyenda
ax.legend(loc='best')

# Guardar
plt.savefig('grafico.png', dpi=300, bbox_inches='tight')
plt.show()
\end{lstlisting}

\subsection{Seaborn: Visualización estadística}

Seaborn está construido sobre matplotlib y facilita la creación de gráficos estadísticos atractivos.

\subsubsection{Configuración de estilo}

\begin{lstlisting}[language=Python]
import seaborn as sns

# Estilos disponibles: darkgrid, whitegrid, dark, white, ticks
sns.set_style("whitegrid")

# Contextos: paper, notebook, talk, poster
sns.set_context("notebook")

# Paletas de colores
sns.set_palette("Set2")
\end{lstlisting}

\subsubsection{Gráficos de distribución}

\begin{lstlisting}[language=Python]
# Histograma con KDE
sns.histplot(data=df, x='columna', kde=True, bins=30)

# Distribución por categoría
sns.histplot(data=df, x='valor', hue='categoria', kde=True)

# KDE plot
sns.kdeplot(data=df, x='columna', fill=True)

# Distribución conjunta
sns.jointplot(data=df, x='var1', y='var2', kind='scatter')
# kind: 'scatter', 'kde', 'hist', 'hex', 'reg'
\end{lstlisting}

\subsubsection{Gráficos categóricos}

\begin{lstlisting}[language=Python]
# Boxplot
sns.boxplot(data=df, x='categoria', y='valor')

# Violin plot
sns.violinplot(data=df, x='categoria', y='valor')

# Swarm plot (puntos sin superposición)
sns.swarmplot(data=df, x='categoria', y='valor')

# Bar plot (con intervalo de confianza)
sns.barplot(data=df, x='categoria', y='valor')

# Count plot (frecuencias)
sns.countplot(data=df, x='categoria')
\end{lstlisting}

\subsubsection{Gráficos de relación}

\begin{lstlisting}[language=Python]
# Scatterplot
sns.scatterplot(data=df, x='var1', y='var2', hue='categoria', 
                size='var3', alpha=0.6)

# Line plot
sns.lineplot(data=df, x='tiempo', y='valor', hue='grupo')

# Regression plot
sns.regplot(data=df, x='var1', y='var2')

# LM plot (con facetas)
sns.lmplot(data=df, x='var1', y='var2', hue='categoria', 
           col='otra_categoria')
\end{lstlisting}

\subsubsection{Pair plot}

Matriz de scatterplots para múltiples variables.

\begin{lstlisting}[language=Python]
sns.pairplot(df, hue='categoria', diag_kind='kde')
plt.show()
\end{lstlisting}

\subsubsection{Heatmap}

\begin{lstlisting}[language=Python]
# Matriz de correlación
corr = df.corr()
sns.heatmap(corr, annot=True, cmap='coolwarm', center=0, 
            fmt='.2f', square=True, linewidths=1)
plt.title('Matriz de Correlación')
plt.show()
\end{lstlisting}

\subsubsection{FacetGrid: Gráficos múltiples}

\begin{lstlisting}[language=Python]
# Grid por categorías
g = sns.FacetGrid(df, col='categoria1', row='categoria2', height=4)
g.map(sns.scatterplot, 'var1', 'var2')
g.add_legend()
plt.show()
\end{lstlisting}

\section{Clase 13: Visualización con R (ggplot2)}

\subsection{Gramática de gráficos}

ggplot2 implementa la "gramática de gráficos" de Leland Wilkinson:

\begin{itemize}
    \item \textbf{Data:} Los datos a visualizar
    \item \textbf{Aesthetics (aes):} Mapeo de variables a propiedades visuales (x, y, color, size, etc.)
    \item \textbf{Geometries (geom):} Tipo de gráfico (puntos, líneas, barras, etc.)
    \item \textbf{Facets:} Paneles múltiples
    \item \textbf{Statistics:} Transformaciones estadísticas
    \item \textbf{Coordinates:} Sistema de coordenadas
    \item \textbf{Themes:} Apariencia general
\end{itemize}

\subsection{Estructura básica}

\begin{lstlisting}[language=R]
library(ggplot2)

ggplot(data = df, aes(x = var1, y = var2)) +
  geom_point() +
  labs(x = "Variable 1", y = "Variable 2", title = "Título")
\end{lstlisting}

El operador \texttt{+} agrega capas al gráfico.

\subsection{Aesthetics (aes)}

\begin{lstlisting}[language=R]
ggplot(df, aes(x = var1, y = var2, 
               color = categoria,    # Color por categoría
               size = var3,          # Tamaño por variable
               shape = grupo,        # Forma por grupo
               alpha = 0.6)) +       # Transparencia
  geom_point()
\end{lstlisting}

\subsection{Geometrías principales}

\subsubsection{Scatterplots}

\begin{lstlisting}[language=R]
# Básico
ggplot(df, aes(x = var1, y = var2)) +
  geom_point()

# Con color y tamaño
ggplot(df, aes(x = var1, y = var2, color = categoria, size = var3)) +
  geom_point(alpha = 0.6) +
  scale_color_brewer(palette = "Set2")
\end{lstlisting}

\subsubsection{Line plots}

\begin{lstlisting}[language=R]
ggplot(df, aes(x = tiempo, y = valor, color = grupo)) +
  geom_line(size = 1) +
  geom_point(size = 2)
\end{lstlisting}

\subsubsection{Histogramas}

\begin{lstlisting}[language=R]
# Básico
ggplot(df, aes(x = columna)) +
  geom_histogram(bins = 30, fill = "steelblue", color = "black")

# Por categoría
ggplot(df, aes(x = valor, fill = categoria)) +
  geom_histogram(bins = 30, alpha = 0.6, position = "identity")
\end{lstlisting}

\subsubsection{Densidad}

\begin{lstlisting}[language=R]
ggplot(df, aes(x = columna, fill = categoria)) +
  geom_density(alpha = 0.5)
\end{lstlisting}

\subsubsection{Boxplots}

\begin{lstlisting}[language=R]
ggplot(df, aes(x = categoria, y = valor, fill = categoria)) +
  geom_boxplot() +
  theme(legend.position = "none")
\end{lstlisting}

\subsubsection{Violin plots}

\begin{lstlisting}[language=R]
ggplot(df, aes(x = categoria, y = valor, fill = categoria)) +
  geom_violin() +
  geom_boxplot(width = 0.1, fill = "white")  # Boxplot superpuesto
\end{lstlisting}

\subsubsection{Bar plots}

\begin{lstlisting}[language=R]
# Frecuencias
ggplot(df, aes(x = categoria)) +
  geom_bar(fill = "steelblue")

# Valores agregados
ggplot(df, aes(x = categoria, y = valor)) +
  geom_bar(stat = "summary", fun = "mean", fill = "steelblue") +
  labs(y = "Media de Valor")
\end{lstlisting}

\subsection{Facetas}

\begin{lstlisting}[language=R]
# Facetas por columna
ggplot(df, aes(x = var1, y = var2)) +
  geom_point() +
  facet_wrap(~ categoria)

# Facetas en grilla
ggplot(df, aes(x = var1, y = var2)) +
  geom_point() +
  facet_grid(categoria1 ~ categoria2)
\end{lstlisting}

\subsection{Estadísticas}

\begin{lstlisting}[language=R]
# Línea de regresión
ggplot(df, aes(x = var1, y = var2)) +
  geom_point() +
  geom_smooth(method = "lm", se = TRUE)  # se = intervalo de confianza

# Smooth (LOESS)
ggplot(df, aes(x = var1, y = var2)) +
  geom_point() +
  geom_smooth(se = TRUE)
\end{lstlisting}

\subsection{Escalas}

\begin{lstlisting}[language=R]
# Escala de colores
ggplot(df, aes(x = var1, y = var2, color = valor)) +
  geom_point() +
  scale_color_gradient(low = "blue", high = "red")

# Escala logarítmica
ggplot(df, aes(x = var1, y = var2)) +
  geom_point() +
  scale_x_log10() +
  scale_y_log10()

# Límites de ejes
ggplot(df, aes(x = var1, y = var2)) +
  geom_point() +
  xlim(0, 100) +
  ylim(0, 50)
\end{lstlisting}

\subsection{Temas}

\begin{lstlisting}[language=R]
# Temas predefinidos
ggplot(df, aes(x = var1, y = var2)) +
  geom_point() +
  theme_minimal()
# Otros: theme_bw(), theme_classic(), theme_dark(), theme_void()

# Personalización
ggplot(df, aes(x = var1, y = var2)) +
  geom_point() +
  theme(
    plot.title = element_text(size = 16, face = "bold"),
    axis.title = element_text(size = 12),
    axis.text = element_text(size = 10),
    legend.position = "bottom",
    panel.grid.major = element_line(color = "gray80"),
    panel.grid.minor = element_blank()
  )
\end{lstlisting}

\subsection{Guardar gráficos}

\begin{lstlisting}[language=R]
# Crear gráfico
p <- ggplot(df, aes(x = var1, y = var2)) +
  geom_point()

# Guardar
ggsave("grafico.png", plot = p, width = 10, height = 6, dpi = 300)
# Formatos: png, pdf, svg, jpg
\end{lstlisting}

\subsection{Extensiones de ggplot2}

\begin{itemize}
    \item \textbf{ggpubr:} Gráficos listos para publicación
    \item \textbf{patchwork:} Combinar múltiples gráficos
    \item \textbf{gganimate:} Animaciones
    \item \textbf{plotly + ggplotly:} Gráficos interactivos
    \item \textbf{ggridges:} Ridge plots
    \item \textbf{GGally:} Pair plots y más
\end{itemize}

\begin{lstlisting}[language=R]
# Ejemplo con GGally
library(GGally)

ggpairs(df, columns = c('var1', 'var2', 'var3'), 
        mapping = aes(color = categoria))
\end{lstlisting}

\section*{Resumen de la Unidad 2}

\begin{itemize}[leftmargin=*]
    \item La \textbf{estadística descriptiva} resume datos con medidas de tendencia central (media, mediana, moda) y dispersión (rango, IQR, varianza, desviación estándar)
    \item El \textbf{análisis de distribuciones} usa histogramas, densidades, boxplots y Q-Q plots
    \item Las \textbf{pruebas de normalidad} (Shapiro-Wilk) validan supuestos
    \item Los \textbf{outliers} se detectan con métodos IQR, Z-score, MAD, Mahalanobis, Isolation Forest
    \item El \textbf{análisis bivariado} explora relaciones con scatterplots, correlación (Pearson, Spearman), tablas de contingencia
    \item \textbf{Matplotlib} es la librería base de Python; \textbf{seaborn} facilita gráficos estadísticos
    \item \textbf{ggplot2} implementa la gramática de gráficos en R con sintaxis consistente y elegante
\end{itemize}

\section*{Conceptos clave}

\begin{itemize}[leftmargin=*]
    \item Media, mediana, moda
    \item Varianza, desviación estándar, IQR
    \item Asimetría (skewness), curtosis (kurtosis)
    \item Histogramas, densidad, boxplots, violin plots
    \item Q-Q plots, test de Shapiro-Wilk
    \item Transformaciones: log, sqrt, Box-Cox
    \item Outliers: IQR, Z-score, MAD, Mahalanobis
    \item Scatterplots, correlación (Pearson, Spearman)
    \item Tablas de contingencia, chi-cuadrado
    \item Matplotlib: figure, axes, subplots
    \item Seaborn: histplot, boxplot, scatterplot, heatmap
    \item ggplot2: aes, geom, facet, theme
\end{itemize}

\section*{Próxima unidad}

En la Unidad 3 nos enfocaremos en \textbf{Modelado Estadístico}: regresión lineal simple y múltiple, regularización (Ridge, LASSO), validación cruzada y evaluación de modelos.

\end{document}
